% !TEX root = task-1-7.tex
% Editable task bank. Append new BankTask blocks at the end of the document.
\documentclass[12pt,a4paper]{article}
\usepackage{physics-textbook}
\usepackage{graphicx,needspace}
\usetikzlibrary{decorations.pathmorphing}
\newif\ifShowSolutions
\ifdefined\TutorEdition
  \ShowSolutionstrue
\else
  \ShowSolutionsfalse % Change to \ShowSolutionstrue to display solutions.
\fi
\newcounter{banktask}
\newsavebox{\BankTaskBox}
\newcommand{\BankTaskContent}[3]{%
  \noindent\textbf{\thebanktask.}\quad #2\par
  \ifShowSolutions
    \Needspace{50mm}% Keep the solution heading with its opening text or diagram.
    #3\par
    \ifdefined\TutorEdition
      \ifcsname TutorSource#1\endcsname
        \nopagebreak\smallskip{\small\csname TutorSource#1\endcsname\par}%
      \fi
    \fi
  \fi
}
\newcommand{\BankTask}[3]{%
  \refstepcounter{banktask}\par\medskip
  \sbox{\BankTaskBox}{\begin{minipage}{\linewidth}%
    \BankTaskContent{#1}{#2}{#3}\end{minipage}}%
  \ifdim\dimexpr\ht\BankTaskBox+\dp\BankTaskBox\relax>\dimexpr\textheight-20mm\relax
    \Needspace{35mm}\BankTaskContent{#1}{#2}{#3}%
  \else
    \Needspace{\dimexpr\ht\BankTaskBox+\dp\BankTaskBox\relax}%
    \noindent\usebox{\BankTaskBox}\par
  \fi
  \medskip
}
\begin{document}
\BookSubsection{1.7}{Масса и плотность}


% Source 60DF4D: https://ege.fipi.ru/bank/questions.php?search=1&proj=BA1F39653304A5B041B656915DC36B38&qid=60DF4D
\BankTask{60DF4D}{
Необходимо собрать экспериментальную установку, с помощью которой можно определить плотность алюминия. Для этого школьник взял стакан с водой и алюминиевый шарик с крючком.

Какое дополнительное оборудование необходимо использовать для проведения этого эксперимента?
\begin{enumerate}
\item[1)] электронные весы
\item[2)] мензурка
\item[3)] секундомер
\item[4)] термометр
\item[5)] пружина
\end{enumerate}
}{
\smallskip\noindent\textbf{Решение от GPT}\par
Электронные весы (1) позволяют измерить массу шарика \(m\), а мензурка (2) — его объём по изменению уровня воды: \(V=V_2-V_1\). Плотность алюминия равна \(\rho=m/(V_2-V_1)\). Секундомер (3) и термометр (4) не нужны, а пружина (5) без известной жёсткости не позволяет измерить массу и не заменяет мензурку. Ответ: 1, 2.
}

\end{document}
